% Generated draft: requires mathematical and source review.
\documentclass[11pt]{article}
\usepackage[margin=25mm]{geometry}
\usepackage{fontspec}
\setmainfont{Latin Modern Roman}
\usepackage{amsmath,amssymb,mathtools,unicode-math}
\setmathfont{Latin Modern Math}
\usepackage{enumitem,needspace,xurl,hyperref}
\setlength{\parindent}{0pt}
\setlength{\parskip}{5pt}
\setlength{\emergencystretch}{3em}
\newcommand{\N}{\mathbb{N}}
\newcommand{\Z}{\mathbb{Z}}
\newcommand{\Q}{\mathbb{Q}}
\newcommand{\R}{\mathbb{R}}
\newcommand{\C}{\mathbb{C}}
\newcommand{\F}{\mathbb{F}}
\newcommand{\A}{\mathbb{A}}
\newcommand{\PP}{\mathbb{P}}
\newcommand{\E}{\mathbb{E}}
\newcommand{\eps}{\varepsilon}
\newcommand{\dd}{\,\mathrm{d}}
\begin{document}
\section*{UnsolvedMath ID: 30004670}
\textbf{Source problem code:} catalogue entry
\subsection{Definitions and mathematical statement}
For sets \(X\) and \(Y\), a multifunction may be described by a relation \(F\subseteq X\times Y\), with values and domain
\[
F(x)=\{y\in Y:(x,y)\in F\},\qquad
\operatorname{dom}(F)=\{x\in X:F(x)\ne\varnothing\}.
\]
This notation permits partial multifunctions; it does not assert that the source intends to include them. A total multifunction would additionally satisfy \(\operatorname{dom}(F)=X\).

\textbf{Requested characterization.} For arbitrary spaces \(X\) and \(Y\) that are effectively Hausdorff in the sense intended by the source, find and prove necessary and sufficient conditions for a multifunction \(F:X\rightrightarrows Y\) to be computable.

\textbf{Unresolved qualifications.} The supplied record does not define an effectively Hausdorff space, specify the effective presentations or representations of \(X\) and \(Y\), or define computability for multifunctions. In particular, it does not specify whether computability concerns producing an allowed output or computing the entire output set in some representation. It also does not specify whether multifunctions must be total or whether their values have additional restrictions. Consequently, the requested characterization is a source-level problem statement whose computational meaning and quantifier scope require clarification; the relational notation above does not resolve these missing conventions.
\subsection{Short English statement}
Characterize computable multifunctions between arbitrary effectively Hausdorff spaces and prove the characterization in both directions. The source does not supply the effective structures, the definition of multifunction computability, or the totality and value restrictions needed for a fully precise formulation. Its open status is supplied dated curation, not an independently verified conclusion.
\subsection{Sources}
\begin{itemize}
\item Supplied local UnsolvedMath registry record, permanent ID \(30004670\), title \emph{Computable Multifunctions Between Effectively Hausdorff Spaces}, source problem code \texttt{catalogue entry}, in the Oberwolfach Reports Open Problems collection. The record gives proposed year 2021 and names no proposer.
\item Supplied source citation: \emph{Computability Theory (2022)}, Oberwolfach Reports, DOI 10.4171/owr/2021/21, pp.\ 1149--1190. \url{https://doi.org/10.4171/owr/2021/21}. This attribution is reproduced from the record; the source has not been independently consulted.
\item Supplied literature source: \emph{Oberwolfach Report 21/2021}. \url{https://ems.press/content/serial-article-files/46899}.
\item Supplied literature source: Schr\"oder, \emph{Admissibly represented spaces and QCB-spaces}. \url{https://arxiv.org/abs/2004.09450}.
\item Literature assessment checked at \texttt{2026-08-22}, according to the supplied record. The supplied assessment labels the problem open, attributes the introduction of effectively Hausdorff spaces and the characterization question to the report, and states that later surveys develop the framework without supplying the requested general characterization. These are supplied dated curation claims, not independently verified facts; time-sensitive status and cited claims require independent verification.
\end{itemize}
\end{document}
